\chapter{ГЕНЕРАЦИЯ КИНЕМАТИЧЕСКОЙ СХЕМЫ НА ОСНОВЕ ФОРМАТА STEP}
\label{cha:step}

\section{Постановка задачи}

Конструкторская документация на робототехнические изделия чаще всего
поставляется в нейтральном формате STEP (ISO~10303)~\cite{iso10303}, который
описывает геометрию сборки и иерархию деталей, но \emph{не содержит} сведений о
подвижных соединениях (сопряжениях, mate-связях). Для использования такой
сборки в симуляции необходимо восстановить её кинематическую схему: дерево
звеньев и суставов с их типами и осями, инерционные характеристики звеньев и
описание модели в формате, пригодном для симулятора. Задача решена в
репозитории \texttt{kinechain} (язык Python, геометрическое ядро
OpenCASCADE~\cite{opencascade} через пакет \texttt{cadquery-ocp}); работа по
данному пункту календарного плана завершена.

\section{Конвейер обработки}

Инструмент организован как линейный конвейер, по одному модулю на стадию:

\begin{enumerate}
  \item \textbf{Загрузка STEP} (\texttt{io\_step}): чтение сборки с сохранением
        иерархии и мировых преобразований; на выходе --- плоский список
        именованных деталей, геометрия которых уже приведена к мировой системе
        координат.
  \item \textbf{Извлечение аналитических поверхностей} (\texttt{surfaces}): из
        граничного представления (B-Rep) каждой детали извлекаются плоскости
        (точка, нормаль) и цилиндры (ось, точка на оси, радиус, осевая
        протяжённость).
  \item \textbf{Обнаружение контактов} (\texttt{contact}): широкая фаза по
        перекрытию ограничивающих параллелепипедов (AABB), узкая фаза ---
        проверки геометрической совместности пар поверхностей; на выходе ---
        типизированные \emph{признаки контакта} пары деталей.
  \item \textbf{Классификация суставов} (\texttt{joints}): детерминированные
        правила распознают по признакам контакта вращательный
        (\texttt{revolute}), поступательный (\texttt{prismatic}) или
        фиксированный (\texttt{fixed}) сустав.
  \item \textbf{Построение кинематического дерева} (\texttt{graph}): остовное
        дерево обходом в ширину от корня; не вошедшие в остов суставы
        становятся замыкающими ограничениями циклов.
  \item \textbf{Экспорт} (\texttt{export\_mjcf}, \texttt{export\_sdf}): запись
        модели в MJCF (MuJoCo) или SDF из общего промежуточного представления,
        с выгрузкой полигональных сеток деталей в STL.
\end{enumerate}

Все геометрические допуски (линейный, угловой, по радиусу, минимальные площади
контакта) собраны в одной структуре и выводятся из диагонали ограничивающего
параллелепипеда всей сборки, что делает конвейер \emph{масштабно-инвариантным}:
одни и те же правила работают на сборках от часового до строительного масштаба
без перенастройки.

\section{Обнаружение контактов}

Широкая фаза отбрасывает заведомо не соприкасающиеся пары деталей по перекрытию
AABB (с зазором, равным линейному допуску). Для оставшихся пар узкая фаза
проверяет два вида совместности поверхностей.

\textbf{Соосность цилиндров.} Для цилиндров с осями $\mathbf{a}_1,\mathbf{a}_2$
и точками на осях $\mathbf{p}_1,\mathbf{p}_2$ требуется одновременное
выполнение условий параллельности осей и совпадения прямых:
\begin{equation}
  \lVert \mathbf{a}_1 \times \mathbf{a}_2 \rVert < \varepsilon_{\angle},
  \qquad
  \bigl\lVert (\mathbf{p}_2 - \mathbf{p}_1) \times \mathbf{a}_1 \bigr\rVert
  < \varepsilon_{d},
  \label{eq:coaxial}
\end{equation}
где $\varepsilon_{\angle}$ и $\varepsilon_{d}$ --- угловой и линейный допуски;
дополнительно радиусы должны совпадать в пределах допуска. Для каждой соосной
пары вычисляется \emph{осевое перекрытие} --- длина общего участка вдоль оси,
используемая затем классификатором.

\textbf{Прилегание плоскостей.} Две плоские грани считаются совмещёнными при
параллельности нормалей и расстоянии между плоскостями в пределах допуска;
признак \emph{встречности} нормалей отличает настоящий контакт <<грань к
грани>> от заподлицо расположенных граней. Площадь контакта вычисляется
\emph{точно} (с точностью до триангуляции): триангуляции обеих граней
проецируются на общую плоскость, и суммируется площадь попарных пересечений
треугольников, найденных алгоритмом отсечения выпуклых многоугольников
Сазерленда--Ходжмана. Это корректно даёт нулевую площадь для компланарных, но
не перекрывающихся граней.

\section{Классификация суставов}

Классификация построена на детерминированных правилах, применяемых в порядке
от наиболее ограничивающего типа к наименее: фиксированный $\rightarrow$
поступательный $\rightarrow$ вращательный (листинг~\ref{lst:joint-rules}).
Правила \emph{консервативны}: если ни одно не сработало, пара остаётся
несоединённой и попадает в диагностику, а не <<приваривается>> молча.

\begin{pseudocode}{Алгоритм классификации сустава по признакам контакта}{lst:joint-rules}
\pind0 Вход: соосные пары цилиндров $CC$, совмещённые пары плоскостей $PP$,
       допуски\\
\pind0 $\Pi \leftarrow \{\,p \in PP :$ нормали встречные, площадь контакта
       $\ge S_{\min}\,\}$\\
\pind0 $N \leftarrow$ различные (с точностью до знака) нормали пар из $\Pi$\\
\pind0 \textit{// 1. Фиксированный}\\
\pind0 если $\Pi \neq \varnothing$ и (в $CC$ есть <<болтовая>> пара
       (перекрытие $<$ диаметра)\\
\pind2 или $|N| \ge 2$ и у $N$ нет общего перпендикуляра):\\
\pind1 вернуть \texttt{fixed} с якорем на плоскости наибольшей площади\\
\pind0 \textit{// 2. Поступательный}\\
\pind0 $\mathbf{s} \leftarrow$ общий перпендикуляр ко всем нормалям из $N$
       (существует при $|N| \ge 2$)\\
\pind0 если $\mathbf{s}$ существует и ось каждой пары из $CC$ параллельна
       $\mathbf{s}$:\\
\pind1 вернуть \texttt{prismatic} с осью $\mathbf{s}$\\
\pind0 \textit{// 3. Вращательный}\\
\pind0 если в $CC$ есть пара: перекрытие $\ge$ диаметра и в $\Pi$ нет
       плоскости с нормалью вдоль её оси:\\
\pind1 вернуть \texttt{revolute} с осью этой пары\\
\pind0 иначе вернуть <<нет сустава>> \quad\textit{// пара остаётся
       несоединённой, выводится в диагностику}
\end{pseudocode}

Содержательная основа правил --- анализ степеней свободы, оставляемых
контактами:
\begin{itemize}
  \item \emph{поступательный}: два и более контактов <<грань к грани>> с
        непараллельными нормалями запрещают все вращения и все перемещения,
        кроме одного --- вдоль общего перпендикуляра $\mathbf{s}$ к нормалям
        контактов (направления скольжения); соосный цилиндр при этом не должен
        противоречить $\mathbf{s}$, иначе скольжение разорвало бы контакт;
  \item \emph{фиксированный}: нормали контактов, не имеющие общего
        перпендикуляра, запрещают и последнее перемещение --- пара сварена
        самой геометрией; короткая <<болтовая>> соосная пара (осевое перекрытие
        меньше диаметра --- шпилька или болт, а не подшипниковая посадка) в
        сочетании с прилеганием плоскостей также трактуется как крепёж;
  \item \emph{вращательный}: длинная соосная пара (перекрытие не меньше
        диаметра) без <<заплечика>> --- прилегающей плоскости с нормалью вдоль
        оси, которая исключила бы свободное вращение.
\end{itemize}

Ключевой фрагмент реализации --- вычисление направления скольжения и цепочка
правил --- приведён в листинге~\ref{lst:kinechain}.

\begin{lstlisting}[language=Python,caption={Направление скольжения и цепочка правил (\texttt{joints.py}, фрагмент)},label=lst:kinechain]
def _slide_direction(normals: list[np.ndarray]) -> np.ndarray | None:
    """Единственное направление перемещения, оставляемое контактами,
    или None, если нормали покрывают все три измерения."""
    if len(normals) < 2:
        return None
    s = np.cross(normals[0], normals[1])
    s = s / np.linalg.norm(s)
    for n in normals[2:]:
        if abs(float(np.dot(n, s))) > _PERP_SIN:
            return None
    return s

def classify(contact: PartContact, tol: Tolerances) -> Joint | None:
    # От наиболее ограничивающего правила к наименее; None — пара
    # остаётся несоединённой и попадает в диагностику.
    return (
        _match_fixed(contact, tol)
        or _match_prismatic(contact, tol)
        or _match_revolute(contact)
    )
\end{lstlisting}

\section{Построение дерева и замыкание циклов}

Найденные суставы образуют неориентированный граф деталей. Корнем по умолчанию
выбирается участвующая в соединениях деталь наибольшего объёма (как правило,
станина или основание). Обходом в ширину от корня строится остовное дерево;
рёбра дерева ориентируются <<родитель $\rightarrow$ потомок>> по порядку
обхода. Суставы, не вошедшие в остов, --- признак замкнутой кинематической
цепи; они выносятся в отдельный список \emph{замыкающих суставов} и
экспортируются как ограничения замыкания цикла (в MJCF --- элемент
\texttt{equality/connect}). Корректность замыкания проверена на классическом
четырёхзвенном механизме: цикл из четырёх вращательных суставов разбивается на
остов из трёх суставов и одно замыкающее ограничение. Результат работы
инструмента на этой сборке показан на рисунке~\ref{fig:step-kin}: слева ---
исходная STEP-сборка (звенья в трёх слоях по высоте, суставы образованы
парами <<палец --- отверстие>>), справа --- восстановленная кинематическая
схема.

\definecolor{kcGround}{rgb}{0.62,0.62,0.62}
\definecolor{kcCrank}{rgb}{0.85,0.36,0.27}
\definecolor{kcCoupler}{rgb}{0.27,0.55,0.85}
\definecolor{kcRocker}{rgb}{0.33,0.70,0.35}

\begin{figure}[ht]
  \centering
  \begin{tabular}{cc}
  % --- а) исходная сборка: косоугольная проекция, мм -> см, ось Z растянута ---
  \begin{tikzpicture}[x={(0.05cm,0cm)}, y={(0.023cm,0.014cm)}, z={(0cm,0.1cm)},
                      line cap=round]
    % станина (z 0..5), пальцы в A и B
    \draw[line width=0.5cm, kcGround]  (2,0,2.5)  -- (58,0,2.5);
    \fill[kcGround]  (0,0,2.5)  circle (0.32cm);
    \fill[kcGround]  (60,0,2.5) circle (0.32cm);
    % видимые участки пальцев A, B в зазоре (z 5..7)
    \draw[line width=0.16cm, kcGround!80!black] (0,0,5)  -- (0,0,7);
    \draw[line width=0.16cm, kcGround!80!black] (60,0,5) -- (60,0,7);
    % кривошип A-C и коромысло B-D (z 7..12)
    \draw[line width=0.5cm, kcCrank]  (1,4,9.5)  -- (9,36,9.5);
    \fill[kcCrank]  (0,0,9.5)  circle (0.3cm);
    \fill[kcCrank]  (10,40,9.5) circle (0.3cm);
    \draw[line width=0.5cm, kcRocker] (59,4,9.5) -- (51,36,9.5);
    \fill[kcRocker] (60,0,9.5) circle (0.3cm);
    \fill[kcRocker] (50,40,9.5) circle (0.3cm);
    % верхушки пальцев станины над кривошипом/коромыслом (z 12..17)
    \draw[line width=0.16cm, kcGround!80!black] (0,0,12)  -- (0,0,17);
    \fill[kcGround!80!black] (0,0,17)  ellipse [x radius=0.1cm, y radius=0.045cm];
    \draw[line width=0.16cm, kcGround!80!black] (60,0,12) -- (60,0,17);
    \fill[kcGround!80!black] (60,0,17) ellipse [x radius=0.1cm, y radius=0.045cm];
    % видимые участки пальцев C, D в зазоре (z 12..14)
    \draw[line width=0.16cm, kcCrank!75!black]  (10,40,12) -- (10,40,14);
    \draw[line width=0.16cm, kcRocker!75!black] (50,40,12) -- (50,40,14);
    % шатун C-D (z 14..19)
    \draw[line width=0.5cm, kcCoupler] (14,40,16.5) -- (46,40,16.5);
    \fill[kcCoupler] (10,40,16.5) circle (0.3cm);
    \fill[kcCoupler] (50,40,16.5) circle (0.3cm);
    % верхушки пальцев C, D над шатуном (z 19..24)
    \draw[line width=0.16cm, kcCrank!75!black]  (10,40,19) -- (10,40,24);
    \fill[kcCrank!75!black]  (10,40,24) ellipse [x radius=0.1cm, y radius=0.045cm];
    \draw[line width=0.16cm, kcRocker!75!black] (50,40,19) -- (50,40,24);
    \fill[kcRocker!75!black] (50,40,24) ellipse [x radius=0.1cm, y radius=0.045cm];
    % подписи суставов и звеньев
    \node[font=\scriptsize, below=5pt] at (0,0,0)   {$A$};
    \node[font=\scriptsize, below=5pt] at (60,0,0)  {$B$};
    \node[font=\scriptsize, above=3pt] at (10,40,24) {$C$};
    \node[font=\scriptsize, above=3pt] at (50,40,24) {$D$};
    \node[font=\scriptsize, below=7pt] at (30,0,0)   {станина};
    \node[font=\scriptsize, left=5pt]  at (5,20,9.5) {кривошип};
    \node[font=\scriptsize, right=5pt] at (55,20,9.5){коромысло};
    \node[font=\scriptsize, above=5pt] at (30,40,19) {шатун};
  \end{tikzpicture}
  &
  % --- б) кинематическая схема (kinematikz), тот же план XY ---
  \begin{tikzpicture}[scale=1.0]
    \coordinate (A) at (0,0);
    \coordinate (B) at (3,0);
    \coordinate (C) at (0.5,2);
    \coordinate (D) at (2.5,2);
    % стойки (станина) в A и B
    \pic at (A) {frame pivot triangle=0.5cm};
    \pic at (B) {frame pivot triangle=0.5cm};
    % звенья
    \draw[knLinkStyle] (A) -- (C);
    \draw[knLinkStyle] (C) -- (D);
    \draw[knLinkStyle] (B) -- (D);
    % шарниры C и D (A и B нарисованы стойками)
    \draw[knLinkStyle, fill=white] (C) circle (\pivotRadius);
    \draw[knLinkStyle, fill=white] (D) circle (\pivotRadius);
    % замыкающее ограничение: сустав D вне остовного дерева
    \draw[dashed, thick, red!65!black] (D) circle (0.24cm);
    \node[font=\scriptsize, text=red!65!black, align=left, anchor=south west]
      at ($(D)+(0.30,0.34)$) {замыкающее\\ограничение};
    % подписи суставов и звеньев
    \node[font=\scriptsize, below left=2pt and 0pt]  at (A) {$A$};
    \node[font=\scriptsize, below right=2pt and 0pt] at (B) {$B$};
    \node[font=\scriptsize, above left=2pt and 0pt]  at (C) {$C$};
    \node[font=\scriptsize, anchor=east] at ($(D)+(-0.28,-0.10)$) {$D$};
    \node[font=\scriptsize, rotate=76, above=2pt] at ($(A)!0.55!(C)$) {кривошип};
    \node[font=\scriptsize, above=2pt] at ($(C)!0.45!(D)$) {шатун};
    \node[font=\scriptsize, rotate=-76, below=2pt] at ($(B)!0.45!(D)$) {коромысло};
  \end{tikzpicture}
  \\[2pt]
  а) исходная STEP-сборка & б) восстановленная кинематическая схема \\
  \end{tabular}
  \caption{Результат работы инструмента на четырёхзвенном механизме: четыре
    вращательных сустава, из которых три образуют остовное дерево и один ---
    замыкающее ограничение цикла}
  \label{fig:step-kin}
\end{figure}

\section{Инерционные характеристики и экспорт}

Масса, центр масс и полный тензор инерции каждого звена вычисляются
\emph{точным} объёмным интегрированием граничного представления средствами
OpenCASCADE (а не по аппроксимирующей сетке): тензор берётся относительно
центра масс, применяется плотность по умолчанию и перевод единиц
мм$\,\rightarrow\,$м. Это устраняет систематическую погрешность инерции,
свойственную интегрированию по триангуляции.

Оба экспортёра используют общее промежуточное представление (дерево, суставы,
замыкающие ограничения) и общую выгрузку сеток STL, отличаясь семантикой
формата: MJCF вкладывает тела друг в друга (фиксированный сустав --- отсутствие
элемента \texttt{joint}), SDF описывает модель плоским списком с явными
суставами и приваркой корня к псевдозвену \texttt{world}. Полученная модель
напрямую загружается подсистемой сцен \texttt{modeuler} (раздел~\ref{cha:sdf})
и физическим ядром.

\section{Тестирование}

Тестовые сборки строятся программно (пакет \texttt{cadquery}) и экспортируются
в STEP непосредственно в ходе теста, поэтому бинарные файлы в репозитории не
хранятся. Геометрия каждой сборки намеренно подобрана так, чтобы разделять
правила классификации:
\begin{itemize}
  \item \emph{шарнир} --- палец выступает за обе щеки кронштейна (нет контакта
        <<заплечиком>>), допустим только \texttt{revolute};
  \item \emph{скреплённые болтами пластины} --- соосные пары <<шпилька ---
        отверстие>> выглядят как вращательная пара, но обязаны
        классифицироваться как \texttt{fixed} (болтовая эвристика);
  \item \emph{ползун на V-образной направляющей} --- только плоскостные
        контакты, общий перпендикуляр нормалей задаёт ось \texttt{prismatic};
  \item \emph{четырёхзвенный механизм} (рисунок~\ref{fig:step-kin}) ---
        замкнутая цепь из четырёх \texttt{revolute}, проверяющая разбиение
        <<остов + замыкающее ограничение>>.
\end{itemize}
Помимо модульных тестов геометрических предикатов и правил, выполняется
сквозная проверка: экспортированная модель MJCF загружается в симулятор MuJoCo
и прогоняется на несколько шагов (тест автоматически пропускается, если MuJoCo
не установлен). Набор исполняется в непрерывной интеграции на Python 3.11 и
3.12.

\section{Результаты}

Работа календарного плана выполнена полностью: реализован сквозной инструмент
построения кинематической схемы из STEP-сборки без явной информации о
сопряжениях. Поддержаны три типа суставов (вращательный, поступательный,
фиксированный) с консервативной детерминированной классификацией по признакам
контактов поверхностей, замыкание циклов кинематических цепей, точные
инерционные характеристики звеньев из объёмного интегрирования B-Rep и экспорт
в два формата (MJCF и SDF) из общего промежуточного представления.
Масштабно-инвариантные допуски позволяют применять инструмент к сборкам
произвольного размера. Корректность подтверждена набором дискриминирующих
тестовых сборок и сквозной загрузкой экспортированных моделей в симулятор.
